\section*{Bab \ref{ch:b1:earthquake-source} --- Sumber gempa bumi}

\begin{solution}{exo:b1:earthquake-source:1}
(a) $\log_{10}(\num{4e17}) = 17.60$, sehingga $M_w = \tfrac23(17.60 - 9.1) = 5.67$.
(b) $\log_{10}(\num{5.59e22}) = 22.75$ dan $M_w = \tfrac23(22.75 - 9.1) =
9.10$. (c) $M_0 = 10^{1.5\times7.0 + 9.1} = 10^{19.6} = \qty{3.98e19}{N.m}$.
\end{solution}

\begin{solution}{exo:b1:earthquake-source:2}
(a) $T = D/\dot s = \qty{3}{m}/\qty{0.025}{m.yr^{-1}} = \qty{120}{yr}$.
(b) Dengan $c = 0.5$, $T = D/(c\dot s) = \qty{240}{yr}$: rayapan mengambil separuh gerak, dan gempa bumi dengan pergeseran yang sama datang setengah sesering.
\end{solution}

\begin{solution}{exo:b1:earthquake-source:3}
Menurut \cref{prop:b1:earthquake-source:moment-steps}, perbandingan momennya adalah
$10^{1.5\times(7.4 - 6.2)} = 10^{1.8} = 63.1$: dibutuhkan 63 gempa bumi bermagnitudo 6.2 untuk melepaskan momen satu bermagnitudo 7.4.
\end{solution}

\begin{solution}{exo:b1:earthquake-source:4}
Dengan $b = 1$, setiap satuan magnitudo membagi jumlah dengan 10: 10 gempa bumi bermagnitudo 5 atau lebih setahun, 1 bermagnitudo 6 atau lebih, dan 0.1 bermagnitudo 7 atau lebih, yaitu satu setiap 10 tahun rata-rata.
\end{solution}

\begin{solution}{exo:b1:earthquake-source:5}
(a) $A = \num{4.3e5} \times \qty{1.5e4}{m} = \qty{6.45e9}{m^2}$, dan
$M_0 = \num{3e10}\times\num{6.45e9}\times4 = \qty{7.74e20}{N.m}$;
$\log_{10}M_0 = 20.89$ dan $M_w = \tfrac23(20.89 - 9.1) = 7.86$. (b) Kedalaman
\qty{10}{km} membagi luas dan momen dengan 1.5, yang menurunkan magnitudo sebesar
$\tfrac23\log_{10}1.5 = 0.12$, menjadi 7.74. (c) Yang pertama, 7.86, dibulatkan menjadi
7.9 dari katalog; yang kedua sedikit rendah. Magnitudo tidak terlalu peka terhadap kedalaman: faktor 1.5 pada luas hanya 0.12 satuan magnitudo.
\end{solution}

\begin{solution}{exo:b1:earthquake-source:6}
(a) $D = M_0/(\mu A) = \num{5.59e22}/(\num{3e10}\times\num{4e5}\times\num{1.5e5})
= \qty{31.1}{m}$. (b) $\qty{31.1}{m}/\qty{0.083}{m.yr^{-1}} = \qty{374}{yr}$.
(c) Dengan fraksi terkunci 0.8, $31.1/(0.8\times0.083) = \qty{468}{yr}$; dengan
0.5, \qty{749}{yr}. Makin sedikit antarmuka terkunci, makin lama pergeseran satu gempa bumi besar terkumpul; rata-rata pergeseran yang ditemukan di sini dengan persegi panjang \qtyproduct{400 x 150}{km} dan $\mu = \qty{30}{GPa}$ sendiri hanyalah perkiraan, dan pergeseran nyatanya jauh dari seragam.
\end{solution}

\begin{solution}{exo:b1:earthquake-source:7}
(a) \omterm{def:b1:geological-maps:attitude}{Jurusnya} $\phi = 329$, \omterm{def:b1:geological-maps:attitude}{arah kemiringan} $a = 059$, \omterm{def:b1:geological-maps:attitude}{kemiringan} $\delta =
\ang{8}$. Maka $\mathbf s = (\sin329, \cos329, 0) = (-0.5150, 0.8572, 0)$,
$\mathbf u = (-\cos8\sin59, -\cos8\cos59, \sin8) = (-0.8488, -0.5100, 0.1392)$,
$\mathbf n = (\sin8\sin59, \sin8\cos59, \cos8) = (0.1193, 0.0717, 0.9903)$,
dan dengan $r = \ang{110}$,
\[
  \mathbf d = \cos110\,\mathbf s + \sin110\,\mathbf u = (-0.6215, -0.7724, 0.1308) .
\]
(b) \omterm{def:b1:earthquake-source:focal-mechanism}{Bidang bantu} memiliki normal $\mathbf d$ dan pergeserannya $\mathbf n$. Dengan normal ke atas, ia miring $\arccos 0.1308 = \ang{82.5}$ ke arah
$\arctan(-0.6215/-0.7724) = 218.8$, sehingga \omterm{def:b1:geological-maps:attitude}{jurusnya} $128.8$; pergeseran
$\mathbf n$ diproyeksikan pada vektor \omterm{def:b1:geological-maps:attitude}{jurus} dan vektor \omterm{def:b1:geological-maps:attitude}{naik-kemiringannya} memberi \omterm{def:b1:faults-folds-fabrics:fault-elements}{rake}
\ang{87.2}. Dibulatkan: $129/83/87$, bidang kedua dari solusi yang dipublikasikan.
(c) $\mathbf P = (\mathbf n - \mathbf d)/\sqrt2$ berarah $221$ dan menunjam
\ang{37}; $\mathbf T = (\mathbf n + \mathbf d)/\sqrt2$ berarah $036$ dan menunjam
\ang{52} (dipublikasikan: $222/38$ dan $36/52$). (d) Bidang dangkal, yang miring
\ang{8} ke arah $059$, sepanjang \omterm{def:b1:active-margins:slab}{slab yang menunjam}: \omterm{def:b1:faults-folds-fabrics:fault-types}{sesar naik} dengan \omterm{def:b1:faults-folds-fabrics:fault-elements}{rake}
\ang{110}, dan, karena miring kurang dari \ang{45}, sebuah \omterm{def:b1:faults-folds-fabrics:fault-types}{sesar dorong}, antarmuka lempeng dari zona subduksi.
\end{solution}

\begin{solution}{exo:b1:earthquake-source:8}
(a) Bidang itu miring ke arah $135$, sehingga $\mathbf n = (\sin135, \cos135, 0) =
(0.707, -0.707, 0)$; pergeseran \omterm{def:b1:faults-folds-fabrics:fault-types}{dekstral} memiliki $r = \ang{180}$, $\mathbf d = -\mathbf s
= (-0.707, -0.707, 0)$. Ke utara, $\mathbf x = (0, 1, 0)$:
$(\mathbf n\cdot\mathbf x)(\mathbf d\cdot\mathbf x) = (-0.707)(-0.707) > 0$, sebuah
dorongan. Ke timur: $(0.707)(-0.707) < 0$, sebuah tarikan. Ke selatan: dorongan; ke barat: tarikan.
(b) $\mathbf P = (\mathbf n - \mathbf d)/\sqrt2 = (1, 0, 0)$: \omterm{def:b1:earthquake-source:focal-mechanism}{sumbu P} horizontal dan berarah $090$--$270$; $\mathbf T = (0, -1, 0)$ berarah
$000$--$180$. (c) Bidang vertikal yang berjurus $135$--$315$ memiliki normal
$\mathbf d$. Dengan pergeseran $\mathbf n$, blok di sisi barat daya, ke arah mana $\mathbf d$ menunjuk,
bergerak ke arah $135$; dilihat dari blok timur laut, menghadap barat daya, itu ke kiri:
pergeseran \omterm{def:b1:faults-folds-fabrics:fault-types}{sinistral}. Menurut
\cref{prop:b1:earthquake-source:auxiliary}, pergeseran ini memberi \omterm{def:b1:earthquake-source:focal-mechanism}{gerak pertama} yang sama di setiap arah.
\end{solution}

\begin{solution}{exo:b1:earthquake-source:9}
(a) $V = \qty{40}{mm.yr^{-1}}$ dan $v = V/4$ pada \qty{10}{km}, sehingga menurut
\cref{prop:b1:earthquake-source:profile}(b) $D = \qty{10}{km}$; langsung,
$D = x/\tan(\pi v/V) = 10/\tan(\pi/4) = \qty{10}{km}$. (b) $v(30) =
(40/\pi)\arctan 3 = \qty{15.9}{mm.yr^{-1}}$. (c) $V/(\pi D) =
\num{0.040}/(\pi\times\num{10000}) = \num{1.27e-6}$ per tahun.
\end{solution}

\begin{solution}{exo:b1:earthquake-source:10}
(a) Momen bin masih $C q^k$ dengan $q = 10^{1.5 - b}$, tetapi kini $q < 1$: momen bin menurun ke arah magnitudo besar, dan bin terbawah, $k = 0$, membawa bagian $C/S = (1 - q)/(1 - q^{K+1}) > 1 - q$, bayangan cermin dari \cref{prop:b1:earthquake-source:largest}. Gempa bumi terkecil yang dihitung kemudian akan mendominasi, dan totalnya akan tumbuh tanpa batas ketika yang makin kecil disertakan.
(b) $1 - 10^{-0.7} = 0.80$: paling sedikit \qty{80}{\percent}.
(c) $\log_{10}N = 8.20 - 1.01\times9 = -0.89$, sehingga $N = 0.13$ per tahun dan sekitar 6.4 dalam lima puluh tahun, tiga kali dua yang teramati. Garis itu, yang dicocokkan antara magnitudo 6 dan 7.5, melebihkan gempa bumi terbesar; entah hukum itu melengkung ke bawah di puncak, seperti yang disarankan ukuran sesar yang terbatas, atau lima puluh tahun terlalu sedikit untuk mencicipi peristiwa yang datang sekali dalam beberapa dekade. Keduanya mungkin benar dengan dua peristiwa.
\end{solution}

\begin{solution}{exo:b1:earthquake-source:11}
(a) Dalam 150 tahun sebuah tiang di $x$ bergerak sebesar $150\,v(x) = 150\times(34/\pi)
\arctan(x/15)$ milimeter relatif terhadap tiang di sesar: \qty{0.11}{m} pada
\qty{1}{km}, \qty{1.28}{m} pada \qty{15}{km}, \qty{2.08}{m} pada \qty{50}{km},
ke selatan untuk sesar \omterm{def:b1:faults-folds-fabrics:fault-types}{dekstral} yang dilihat dari barat. (b) Defisit pergeseran adalah
$s = \qty{0.034}{m.yr^{-1}}\times\qty{150}{yr} = \qty{5.1}{m}$. Selama gempa bumi tiang-tiang bergerak sebesar $(s/\pi)\arctan(D/x)$: \qty{2.44}{m} pada
\qty{1}{km}, \qty{1.28}{m} pada \qty{15}{km}, \qty{0.47}{m} pada \qty{50}{km}.
(c) Jumlahnya adalah $0.11 + 2.44$, $1.28 + 1.28$ dan $2.08 + 0.47$, semuanya
\qty{2.55}{m} $= s/2$: setelah gempa bumi pagar lurus kembali di sisi timur, bergeser sebesar $s/2$, dan sisi barat, menurut simetri, sebesar $-s/2$.
\end{solution}

\begin{solution}{exo:b1:earthquake-source:12}
(a) $\mathbf P = (0, 1, 0)$ dan $\mathbf T = (1, 0, 0)$. Pasangan
$\mathbf n = (\mathbf T + \mathbf P)/\sqrt2 = (0.707, 0.707, 0)$,
$\mathbf d = (\mathbf T - \mathbf P)/\sqrt2 = (0.707, -0.707, 0)$ memenuhi
$\mathbf P = (\mathbf n - \mathbf d)/\sqrt2$ dan $\mathbf T = (\mathbf n +
\mathbf d)/\sqrt2$. Kedua vektor horizontal, sehingga kedua \omterm{def:b1:earthquake-source:focal-mechanism}{bidang nodal} vertikal: bidang dengan normal $\mathbf n$ (ke arah $045$) berjurus $135$--$315$, bidang dengan normal $\mathbf d$ (ke arah $135$) berjurus $045$--$225$. (b) Pada bidang yang berjurus $315$, blok di sisi timur laut bergerak sepanjang $\mathbf d$, ke arah $135$: dilihat dari blok lain, ke kanan, pergeseran \omterm{def:b1:faults-folds-fabrics:fault-types}{dekstral}. Pada bidang yang berjurus $045$, blok di sisi tenggara bergerak sepanjang $\mathbf n$, ke arah $045$: pergeseran \omterm{def:b1:faults-folds-fabrics:fault-types}{sinistral}. (c) Sesar baru membentuk $\ang{45} - \phi/2 =
\ang{30}$ dengan $\sigma_1$: dengan $\sigma_1$ berarah $000$, sesar akan berjurus $030$ (\omterm{def:b1:faults-folds-fabrics:fault-types}{sinistral}) dan $330$ (\omterm{def:b1:faults-folds-fabrics:fault-types}{dekstral}), \ang{15} dari \omterm{def:b1:earthquake-source:focal-mechanism}{bidang nodal}. \omterm{def:b1:earthquake-source:focal-mechanism}{Sumbu P} ditetapkan pada \ang{45} dari \omterm{def:b1:earthquake-source:focal-mechanism}{bidang nodal} oleh geometri, apa pun gesekan batuannya, sehingga ia berimpit dengan $\sigma_1$ hanya jika sesar tergelincir pada \ang{45} darinya. (d) Sesar yang berjurus $040$ berjarak \ang{5} dari \omterm{def:b1:earthquake-source:focal-mechanism}{bidang nodal} yang berjurus $045$: bidang itu kemungkinan adalah sesarnya, dan pergeserannya \omterm{def:b1:faults-folds-fabrics:fault-types}{sinistral}.
\end{solution}

\begin{solution}{pb:b1:earthquake-source:1}
\textbf{1.} Interval-interval itu adalah $1290 - 1050 = 240$, $1480 - 1290 = 190$ dan
$1712 - 1480 = 232$ tahun; rata-ratanya adalah \qty{220.7}{yr}.

\textbf{2.} Rata-rata pergeseran adalah $(7.0 + 6.0 + 5.5 + 7.5)/4 = \qty{6.5}{m}$, dan
$T = 6.5/0.030 = \qty{216.7}{yr}$, dekat dengan rata-rata interval.

\textbf{3.} Pergeseran ketiga gempa bumi terakhir, $6.0 + 5.5 + 7.5 =
\qty{19.0}{m}$, terkumpul antara 1050 dan 1712, dalam 662 tahun: $\qty{19.0}{m}
/\qty{662}{yr} = \qty{28.7}{mm.yr^{-1}}$, \qty{4}{\percent} lebih kecil daripada
geodetik \qty{30}{mm.yr^{-1}}. Keduanya sesuai dalam ketidakpastian tanggal dan pergeseran: segmen itu melepaskan hampir seluruh pembebanannya dalam gempa-gempa besarnya.

\textbf{4.} $2025 - 1712 = 313$ tahun, dan defisitnya adalah
$\qty{0.030}{m.yr^{-1}}\times\qty{313}{yr} = \qty{9.39}{m}$.

\textbf{5.} Defisit melampaui setiap pergeseran di parit, yang terbesar \qty{7.5}{m}, dan waktu sejak 1712 melampaui interval terpanjang, 240 tahun. Dalam arti itu segmen tersebut sudah terlambat; itu tidak berarti gempa bumi akan segera terjadi, hanya bahwa sesar itu sudah memegang lebih banyak pergeseran daripada yang pernah dilepaskannya sekaligus dalam catatan.

\textbf{6.} Interval berkisar dari 190 hingga 240 tahun dan pergeserannya dari \qtyrange{5.5}{7.5}{m}: siklusnya tidak periodik, dan kekuatan sesar pada gempa berikutnya tidak diketahui. Parit memberi rata-rata dan sebaran, dari mana sebuah probabilitas dapat diperkirakan, bukan tanggal.

\textbf{7.} $A = \qty{250}{km}\times\qty{15}{km} = \qty{3750}{km^2} =
\qty{3.75e9}{m^2}$.

\textbf{8.} $M_0 = \num{3e10}\times\num{3.75e9}\times\num{9.39} =
\qty{1.06e21}{N.m}$.

\textbf{9.} $M_0 = \num{3e10}\times\num{3.75e9}\times\num{7.5} =
\qty{8.44e20}{N.m}$, sehingga $M_w = \tfrac23(20.926 - 9.1) = 7.88$.

\textbf{10.} Rasionya adalah $9.39/7.5 = 1.25$, yaitu $\tfrac23\log_{10}1.25 =
0.07$ satuan magnitudo.

\textbf{11.} Setengah luas menggandakan momen dan menurunkan magnitudo sebesar
$\tfrac23\log_{10}2 = 0.20$: $M_w = 7.75$.

\textbf{12.} $\mu A V = \num{3e10}\times\num{3.75e9}\times\num{0.030} =
\qty{3.38e18}{N.m}$ per tahun.

\textbf{13.} $b = 1$. $N(5) = 10^{4.3 - 5} = 0.20$ per tahun dan $N(6) = 0.020$ per
tahun.

\textbf{14.} $N(7) = 0.0020$ per tahun: satu setiap 500 tahun; $N(8) = 0.00020$ per
tahun: satu setiap \num{5000} tahun.

\textbf{15.} Hukum memberi satu gempa bumi bermagnitudo 7.95 atau lebih setiap
$1/10^{4.3 - 7.95} = \num{4500}$ tahun, parit satu setiap 220 tahun, dua puluh kali lebih
sering. Hukum dicocokkan dengan banyak gempa bumi kecil dari semua sesar di wilayah itu;
diekstrapolasi ke gempa bumi besar dari satu segmen, ia meremehkannya parah. Perulangan
gempa bumi besar pada sesar diambil dari catatannya sendiri dan dari pembebanannya, bukan
dari yang kecil.

\textbf{16.} $\num{5.5e17}/\num{3.38e18} = 0.16$: gempa bumi sedang melepaskan
\qty{16}{\percent} dari momen yang dikumpulkan segmen itu.

\textbf{17.} Menurut \cref{prop:b1:earthquake-source:moment-steps},
$10^{1.5\times(7.95 - 6)} = 10^{2.92} \approx 840$ kali lebih banyak.

\textbf{18.} Dengan $b = 1$, momen gempa bumi didominasi oleh yang terbesar: apa pun yang
dilakukan gempa sedang, mereka melepaskan bagian kecil, di sini \qty{16}{\percent}, dan sisanya
\qty{84}{\percent}, defisit pergeseran, menunggu gempa bumi besar.

\textbf{19.} $v(-x) = -v(x)$ untuk setiap pasangan, sehingga profilnya antisimetris.
Jauh dari sesar kecepatannya menuju $\pm V/2 = \pm\qty{15}{mm.yr^{-1}}$.

\textbf{20.} $D = x/\tan(\pi v/V) = 15/\tan(\pi\times7.5/30) = 15/\tan\ang{45} =
\qty{15}{km}$: pada \omterm{def:b1:earthquake-source:locking-depth}{kedalaman penguncian} kecepatannya adalah $V/4$.

\textbf{21.} Pada \qty{5}{km}: $D = 5/\tan(\pi\times3.1/30) = 5/\tan\ang{18.6} =
\qty{14.9}{km}$; pada \qty{40}{km}: $D = 40/\tan\ang{69.6} = \qty{14.9}{km}$. Kecepatannya
diberikan hingga \qty{0.1}{mm.yr^{-1}}, dan pembulatan itu saja menggeser $D$ sekitar
\qty{0.1}{km}; dengan data sungguhan, derau tiap stasiun akan lebih besar, dan \omterm{def:b1:earthquake-source:locking-depth}{kedalaman
penguncian} akan dicocokkan dengan semua stasiun sekaligus.

\textbf{22.} \omterm{def:b1:earthquake-source:locking-depth}{Kedalaman penguncian}, \qty{15}{km}, sama dengan kedalaman gempa bumi terdalam. Bagian yang terkunci dari sesar adalah bagian rapuhnya, \omterm{def:b1:stress-strain-rheology:strength-envelope}{lapisan seismogenik}, tempat gesekan menahannya di antara gempa bumi dan ia gagal di dalamnya; di bawahnya, batuan yang lebih hangat merayap (\cref{ch:b1:stress-strain-rheology}) dan sesar itu tergelincir mantap, yaitu pergeseran pada laju $V$ di bawah $D$ dalam model.

\textbf{23.} $V/(\pi D) = \num{0.030}/(\pi\times\num{15000}) = \num{6.4e-7}$
per tahun.

\textbf{24.} Dengan $M_0 = \qty{1.056e21}{N.m}$, $\log_{10}M_0 = 21.024$ dan
$M_w = \tfrac23(21.024 - 9.1) = 7.949$. Hingga 0.01, gempa bumi yang akan melepaskan seluruh
defisit segmen itu memiliki \omterm{def:b1:earthquake-source:mw}{magnitudo momen} \textbf{7.95}.
\end{solution}
